% Vietnamese translation for OI Public Library
% Source-File: IOI中国国家候选队论文/2005/朱晨光/朱晨光.doc
\documentclass[11pt]{article}
\usepackage{oipl-vn}

\newcommand{\Oh}{\mathcal{O}}
\newcommand{\ancestor}{\operatorname{ancestor}}
\newcommand{\lca}{\operatorname{lca}}

\title{Bàn sơ về ứng dụng của tư tưởng nhân đôi trong thi tin học}
\author{Zhu Chenguang\\Trường THPT số 1 Vu Hồ, An Huy}
\date{2005}

\begin{document}
\maketitle

\origtitle{Bàn sơ về ứng dụng của tư tưởng nhân đôi trong thi tin học}
\sourcepath{IOI China candidate team papers / 2005 / Zhu Chenguang.doc}

\renewcommand{\abstractname}{Tóm tắt}
\begin{abstract}
Tư tưởng nhân đôi là một cách nghĩ độc đáo và khéo léo để giải các bài toán
tin học. Bài viết phân tích hai hướng ứng dụng của tư tưởng này trong thi tin
học. Phần đầu nêu vai trò và cách dùng cơ bản của nhân đôi; phần thứ hai trình
bày cách tăng tốc chuyển trạng thái khi quy luật biến đổi không đổi; phần thứ
ba trình bày cách tăng tốc thao tác trên đoạn; phần thứ tư thảo luận vì sao
ta thường mở rộng phạm vi gấp đôi mỗi lần thay vì gấp ba hay nhiều hơn; phần
cuối tổng kết lại ý nghĩa và cách vận dụng linh hoạt tư tưởng nhân đôi.
\end{abstract}

\noindent\textbf{Từ khóa:} tư tưởng nhân đôi; quy luật biến đổi; chuyển trạng
thái; thao tác trên đoạn.

\section{Dẫn nhập}

Nhân đôi là một tư tưởng rất tinh tế và được dùng rộng rãi trong các kỳ thi
tin học. Tuy có thể xuất hiện trong nhiều hoàn cảnh khác nhau, bản chất chung
của nó là: dựa trên thông tin đã tính được, mỗi lần mở rộng phạm vi đang xét
lên gấp đôi để tăng tốc thao tác. Sắp xếp trộn quen thuộc thực ra cũng là một
ứng dụng kinh điển của tư tưởng nhân đôi.

Trong giải bài tập tin học, tư tưởng nhân đôi thường có hai hướng ứng dụng
chính:
\begin{itemize}
  \item tăng tốc chuyển trạng thái khi quy luật biến đổi giống nhau;
  \item tăng tốc thao tác trên đoạn.
\end{itemize}
Hai phần tiếp theo sẽ lần lượt thảo luận hai hướng này.

\section{Ứng dụng thứ nhất: tăng tốc chuyển trạng thái}

Trước hết xét một ví dụ rất đơn giản: biết số thực $a$, tính $a^{17}$.

Cách trực tiếp là đặt $b=a$, rồi lặp $16$ lần phép gán $b=b\cdot a$. Như vậy
cần $16$ phép nhân. Một cách khác là tính
\[
  a_0=a,\qquad a_1=a^2,\qquad a_2=a^4,\qquad
  a_3=a^8,\qquad a_4=a^{16},
\]
trong đó
\[
  a_i=a_{i-1}^2\quad (1\le i\le 4).
\]
Để có $a_0,\ldots,a_4$ chỉ cần $4$ phép nhân, rồi
\[
  a^{17}=a\cdot a^{16}=a_0a_4,
\]
nên tổng cộng chỉ cần $5$ phép nhân.

Mở rộng ra, giả sử cần tính $a^n$. Viết $n$ dưới dạng nhị phân, tức là lấy
những bit khác $0$:
\[
  n=2^{b_1}+2^{b_2}+\cdots+2^{b_w},\qquad
  b_1<b_2<\cdots<b_w.
\]
Từ $a$ ta lần lượt bình phương để có
\[
  a^{2^0},a^{2^1},a^{2^2},\ldots,a^{2^{b_w}}.
\]
Theo quy tắc lũy thừa,
\[
  a^n=a^{2^{b_1}}a^{2^{b_2}}\cdots a^{2^{b_w}}.
\]
Vì biểu diễn nhị phân của $n$ có không quá $\lfloor \log_2 n\rfloor+1$ bit
khác $0$, số phép nhân chỉ còn $\Oh(\log n)$ thay vì $\Oh(n)$. Trong cài đặt,
không cần lưu toàn bộ dãy lũy thừa; khi duyệt các bit từ thấp lên cao, ta vừa
bình phương vừa nhân vào kết quả nếu bit hiện tại bằng $1$.

Thuật toán lũy thừa nhanh giảm số phép nhân vì nó biến việc tính $a^n$ thành
việc nhân một số ít lũy thừa có số mũ là lũy thừa của $2$. Mỗi phần tử sau
trong dãy
\[
  a,a^2,a^4,a^8,\ldots
\]
được tính từ phần tử trước bằng đúng một phép bình phương, trong khi cách thô
cần nhiều phép nhân lặp lại.

Trong thực tế, $a$ có thể là số thực, ma trận, hoặc một trạng thái trừu
tượng; quy luật biến đổi cũng có thể là phép nhân ma trận, chuyển trạng thái
trong quy hoạch động, v.v. Chỉ cần thỏa hai điều kiện sau, ta có thể dùng tư
tưởng nhân đôi:
\begin{enumerate}
  \item quy luật biến đổi ở mỗi bước là giống nhau;
  \item phép ghép các quy luật biến đổi thỏa tính kết hợp.
\end{enumerate}
Trong ví dụ lũy thừa, quy luật là phép nhân, và phép nhân thỏa tính kết hợp.

\subsection{Ví dụ: chuyển động của xúc xắc}

\paragraph{Mô tả bài toán.}
Cho một xúc xắc sáu mặt; mỗi mặt có một trọng số nguyên từ $1$ tới $50$. Xúc
xắc di chuyển trên một dải rộng $4$ hàng và kéo dài vô hạn sang trái, sang
phải. Tọa độ cột là các số nguyên
\[
  \ldots,-3,-2,-1,0,1,2,3,\ldots,
\]
còn tọa độ hàng là $1,2,3,4$. Mỗi ô vừa khít một mặt của xúc xắc. Mỗi bước,
xúc xắc có thể lăn sang trước, sau, trái hoặc phải một ô, nhưng không được ra
khỏi dải. Chi phí của một bước là trọng số của mặt hướng lên sau khi lăn.

Biết vị trí hiện tại và hướng các mặt của xúc xắc, hãy tìm chi phí nhỏ nhất để
di chuyển tới một vị trí mới. Trị tuyệt đối của tọa độ cột có thể tới
$10^9$.

\paragraph{Phân tích.}
Nếu không xét sự chênh lệch cột quá lớn, bài toán có thể giải bằng quy hoạch
động. Mỗi ô được tách thành $24$ trạng thái theo hướng của xúc xắc, nên mỗi
cột có $24\cdot 4=96$ trạng thái. Ta có thể coi bài toán là đồ thị phân tầng:
từ $96$ trạng thái của cột này chuyển sang $96$ trạng thái của cột kế tiếp,
rồi tiếp tục tới cột đích. Mỗi lần chuyển giữa hai cột kề nhau có thể dùng
Dijkstra để tính chi phí nhỏ nhất giữa các trạng thái. Nhưng độ phức tạp tỉ
lệ với độ lệch cột $n$, trong khi $n$ có thể tới $10^9$, nên cách này không
đủ.

Ta kiểm tra hai điều kiện để dùng nhân đôi. Nếu xúc xắc đi từ trạng thái $A$
ở cột $1$ tới trạng thái $B$ ở cột $2$ với chi phí $c$, thì từ trạng thái
tương ứng $A$ ở cột $2$ tới trạng thái $B$ ở cột $3$ cũng có chi phí $c$.
Tổng quát hơn, nếu từ cột $i$ ở trạng thái $A$ tới cột $i+k$ ở trạng thái
$B$ mất chi phí $c$, thì từ cột $j$ ở trạng thái $A$ tới cột $j+k$ ở trạng
thái $B$ cũng mất chi phí $c$. Đây là quy luật biến đổi giống nhau.

Mặt khác, quy hoạch động có tính chất cấu trúc con tối ưu. Nếu đi từ $A$ tới
$C$ qua $B$, thì chi phí nhỏ nhất của toàn đường bằng tổng chi phí nhỏ nhất
của các đoạn tương ứng, và thứ tự ghép các quy luật chuyển có thể xem là kết
hợp. Do đó có thể dùng nhân đôi.

Giả sử đích nằm bên phải điểm xuất phát. Trước hết dùng Dijkstra để tính ma
trận chuyển $M_1$ kích thước $96\times 96$, trong đó $M_1[A,B]$ là chi phí nhỏ
nhất để đi từ trạng thái $A$ của một cột sang trạng thái $B$ của cột ngay bên
phải. Từ $M_1$ tính được $M_2$, chi phí đi xa $2$ cột; rồi $M_4,M_8,\ldots$.
Phép ghép hai ma trận chuyển là phép nhân trong đại số min-plus:
\[
  (X\circ Y)[i,j]=\min_k\{X[i,k]+Y[k,j]\}.
\]

Nếu độ lệch cột là
\[
  n=2^{b_1}+2^{b_2}+\cdots+2^{b_w},
\]
ta xuất phát từ vector trạng thái ban đầu và lần lượt áp dụng các ma trận
$M_{2^{b_1}},M_{2^{b_2}},\ldots,M_{2^{b_w}}$. Kết quả là chi phí nhỏ nhất từ
điểm xuất phát tới mọi trạng thái của cột đích.

Dijkstra ở bước đầu có thể xem là hằng số vì số trạng thái trong một số cột
cần xét là hằng. Mỗi lần nhân đôi ma trận chuyển mất khoảng $96^3$, và số lần
nhân đôi là $\log_2 n$, nên độ phức tạp là
\[
  \Oh(96^3\log n).
\]

Ví dụ này làm rõ một điểm quan trọng: đối tượng mà nhân đôi tác động thực ra
là \term{quy luật biến đổi}, không phải trạng thái cụ thể. Nhân đôi tính ra
``lượng biến đổi'' từ một trạng thái tới một trạng thái khác. Nếu trạng thái
ban đầu là $A$, trạng thái cuối là $B$, và lượng biến đổi là $c$, kết quả có
thể viết trừu tượng là $A*c$, trong đó dấu $*$ biểu thị cách áp dụng quy luật
biến đổi.

Một khuôn giả mã tổng quát là:
\[
\begin{array}{l}
\textbf{DoublingAlgorithm}(A,n):\\
\quad b\gets a,\quad c\gets \text{quy luật đơn vị};\\
\quad \textbf{while }n>0\textbf{ do}\\
\quad\quad \textbf{if }n\text{ lẻ then }c\gets c\circ b;\\
\quad\quad b\gets b\circ b;\\
\quad\quad n\gets \lfloor n/2\rfloor;\\
\quad \textbf{return }A*c.
\end{array}
\]
Ở đây $a,b,c$ đều là quy luật biến đổi; $b\circ b$ biểu thị quy luật tương
đương với việc áp dụng $b$ hai lần liên tiếp.

Cần nhấn mạnh lại: quy luật biến đổi phải giống nhau và phép ghép phải thỏa
tính kết hợp. Phép chia không thỏa điều kiện này nếu giữ nguyên dạng:
\[
  a_1/a_2/a_3/\cdots/a_n
  \ne a_1/(a_2/a_3/\cdots/a_n).
\]
Muốn dùng nhân đôi, ta phải chuyển nó thành một phép hợp lệ, hoặc tìm phương
pháp khác.

\section{Ứng dụng thứ hai: tăng tốc thao tác trên đoạn}

Trong thao tác trên đoạn có nhiều cấu trúc mạnh, tiêu biểu là cây đoạn. Tư
tưởng nhân đôi cũng xử lý được nhiều bài toán đoạn, và trong một số trường
hợp đặc biệt còn rất gọn.

\subsection{Khuôn xử lý đoạn bằng nhân đôi}

\paragraph{Tiền xử lý.}
Với mỗi vị trí $A$, lưu tính chất của các đoạn
\[
  [A,A+2^0-1],\ [A,A+2^1-1],\ [A,A+2^2-1],\ldots,
  [A,A+2^k-1],
\]
trong đó $2^k$ không vượt quá độ dài cần xét. Khi tính tính chất của đoạn
$[A,A+2^{i+1}-1]$, ta ghép hai kết quả đã biết:
\[
  [A,A+2^i-1]\quad\text{và}\quad
  [A+2^i,A+2^{i+1}-1].
\]

\paragraph{Trả lời truy vấn.}
Với một đoạn $[A,B]$, có hai cách lấy kết quả:
\begin{itemize}
  \item Nếu việc chồng lấn đoạn không ảnh hưởng tới kết quả, chẳng hạn lấy
  min hoặc max, đặt
  \[
    k=\lfloor \log_2(B-A+1)\rfloor.
  \]
  Kết quả của $[A,B]$ được ghép từ hai đoạn
  \[
    [A,A+2^k-1]\quad\text{và}\quad [B-2^k+1,B].
  \]
  Hai đoạn có thể chồng lên nhau nhưng không làm sai kết quả. Thời gian trả
  lời là $\Oh(1)$.

  \item Nếu chồng lấn làm thay đổi kết quả, chẳng hạn bài đếm, ta phân rã
  $[A,B]$ thành các đoạn có độ dài là lũy thừa của $2$ theo biểu diễn nhị
  phân. Mỗi lần lấy đoạn lớn nhất có thể, nên số đoạn không quá
  $\lfloor\log_2 N\rfloor+1$. Thời gian trả lời là $\Oh(\log N)$.
\end{itemize}

Tổng quát, tiền xử lý mất $\Oh(N\log N)$ thời gian và bộ nhớ
$\Oh(N\log N)$; mỗi truy vấn mất $\Oh(1)$ hoặc $\Oh(\log N)$ tùy bài.

\subsection{Ví dụ: RMQ tổng quát}

\paragraph{Mô tả bài toán.}
Cho mảng $A$ độ dài $n$. Mỗi truy vấn $(i,j)$ với $i\le j$ yêu cầu tìm chỉ số
của phần tử nhỏ nhất trong $A[i..j]$.

\paragraph{Phân tích.}
Dựng bảng $B$, trong đó $B[i,k]$ là chỉ số của phần tử nhỏ nhất trong đoạn
\[
  A[i..i+2^k-1].
\]
Theo khuôn trên, bảng này được tính trong $\Oh(N\log N)$. Với truy vấn
$(i,j)$, đặt
\[
  k=\lfloor \log_2(j-i+1)\rfloor.
\]
Ta chỉ cần so sánh
\[
  A[B[i,k]]
  \quad\text{và}\quad
  A[B[j-2^k+1,k]]
\]
để lấy chỉ số nhỏ hơn. Vì phép min không bị ảnh hưởng bởi chồng lấn, thời gian
trả lời là $\Oh(1)$. Đây là một ứng dụng rất gọn của nhân đôi trong giảm độ
phức tạp.

\subsection{Ví dụ: tổ tiên chung gần nhất trên cây}

\paragraph{Mô tả bài toán.}
Cho một cây có $n$ đỉnh. Mỗi truy vấn hỏi tổ tiên chung gần nhất của hai đỉnh
$i,j$; số truy vấn là $m$.

\paragraph{Phân tích.}
Nếu mỗi lần tìm trực tiếp, độ phức tạp có thể tới $\Oh(MN)$. Dùng nhân đôi,
với mỗi đỉnh $A$ ta lưu:
\[
  A,\ \text{cha của }A,\ \text{tổ tiên cách }2^1\text{ cạnh},\
  \text{tổ tiên cách }2^2\text{ cạnh},\ldots
\]
Nếu xử lý các đỉnh theo thứ tự duyệt trước, ta có thể dùng kết quả của tổ tiên
đã tính trước để điền bảng. Tiền xử lý mất $\Oh(N\log N)$.

Để trả lời truy vấn $A,B$, mục tiêu là tìm đỉnh cùng xuất hiện trong hai dãy
tổ tiên với độ sâu lớn nhất. Một cách chuẩn là:
\begin{enumerate}
  \item Đưa đỉnh sâu hơn lên cùng độ sâu với đỉnh còn lại bằng các bước nhảy
  $2^k$.
  \item Nếu hai đỉnh đã trùng nhau, đó là LCA.
  \item Ngược lại, từ $k$ lớn xuống nhỏ, nếu tổ tiên $2^k$ của hai đỉnh khác
  nhau thì nâng cả hai lên mức đó.
  \item Cha của hai đỉnh sau bước trên là LCA.
\end{enumerate}
Mỗi truy vấn mất $\Oh(\log N)$.

Bản gốc trình bày một biến thể nhị phân theo dãy tổ tiên của một đỉnh và dùng
bảng nhân đôi để tìm tổ tiên cùng tầng của đỉnh còn lại. Dù cách viết khác
nhau, bản chất vẫn là tiền xử lý các bước nhảy lũy thừa của $2$ và dùng chúng
để thu hẹp khoảng tìm kiếm.

Tài liệu tham khảo của tác giả còn nêu cách chuyển LCA về bài toán RMQ
$\pm1$, có thể đạt tiền xử lý $\Oh(N)$ và trả lời $\Oh(1)$, nhưng cài đặt
phức tạp hơn nên bài viết không triển khai.

Hai ví dụ trên cho thấy nhân đôi có thể xử lý thao tác trên đoạn. Trong RMQ và
LCA, nhân đôi chưa hẳn luôn là phương án mạnh nhất về mặt lý thuyết, nhưng
cài đặt đơn giản, dễ nghĩ và rất hữu dụng. Khi xây dựng suffix array, nhân đôi
còn đóng vai trò cực kỳ quan trọng.

Một điểm đáng chú ý là trong bài LCA và suffix array, đoạn không phải lúc nào
cũng lộ rõ trong đề. Nhiều khi đoạn là cấu trúc do ta tự xây dựng. Điều này
cho thấy nhân đôi là một tư tưởng, không phải một khuôn cứng nhắc; thí sinh
cần dựa vào tư tưởng ấy để tạo ra thuật toán khéo léo và hiệu quả.

\section{Một thảo luận thú vị}

Bản chất của nhân đôi là mỗi lần mở rộng phạm vi lên gấp đôi. Vậy có thể mỗi
lần mở rộng gấp ba, gấp bốn hoặc nhiều hơn không?

Gọi cách mỗi lần mở rộng theo cơ số $k$ là ``$k$-tăng'', với $k\ge 2$. Tư
tưởng nhân đôi thông thường chính là trường hợp $k=2$. Với cơ số $k$, số lần
mở rộng tối đa xấp xỉ $\log_k N$. Nhưng để mở rộng phạm vi lên $k$ lần, nếu
không tiếp tục dùng nhân đôi bên trong, cần khoảng $k-1$ thao tác ghép. Do đó
độ phức tạp tương ứng là
\[
  \Oh((k-1)\log_k N).
\]

Khi $k=2$, giá trị này là $\log_2 N$. So sánh với $k\ge 3$, ta cần xét
\[
  (k-1)\log_k N \quad\text{và}\quad \log_2 N.
\]
Vì
\[
  \log_k N=\frac{\log_2 N}{\log_2 k},
\]
nên tỉ số là
\[
  \frac{(k-1)\log_k N}{\log_2 N}
  =\frac{k-1}{\log_2 k}.
\]
Với $k=2$ tỉ số bằng $1$. Với $k\ge3$, ta có
\[
  k-1>\log_2 k,
\]
nên tỉ số lớn hơn $1$. Nói cách khác,
\[
  (k-1)\log_k N>\log_2 N\qquad (k\ge3).
\]
Điều này giải thích từ góc độ toán học vì sao mở rộng gấp đôi là lựa chọn
hiệu quả tự nhiên.

\section{Kết luận}

Bài viết đã giới thiệu ngắn gọn vai trò của tư tưởng nhân đôi và hai hướng
ứng dụng chính: tăng tốc chuyển trạng thái khi quy luật biến đổi giống nhau,
và tăng tốc thao tác trên đoạn. Nguồn gốc hiệu quả của nhân đôi nằm ở việc tận
dụng đúng mức các kết quả đã tính. Từ phân tích ở trên, có thể thấy nhân đôi
gần như tận dụng việc tái sử dụng này tới cực hạn. Vì vậy các thuật toán được
thiết kế theo tư tưởng nhân đôi thường ít phép tính dư thừa, vừa hiệu quả vừa
gọn.

Nắm vững nhân đôi giúp ta tìm được lối đi riêng khi phân tích bài toán, nhanh
chóng đạt được thuật toán có độ phức tạp thời gian, bộ nhớ và cài đặt tương
đối thấp. Tuy bài viết đã nêu hai dạng ứng dụng và khuôn tổng quát, vai trò
của tư tưởng nhân đôi không dừng ở đó; những biến hóa trong thực tế không thể
gói hết trong một vài mô hình. Hy vọng bạn đọc từ đây cảm nhận được chiều sâu
của tư tưởng nhân đôi, rồi trong quá trình luyện tập dần tích lũy kinh nghiệm,
đưa nó tự nhiên vào cách suy nghĩ của mình.

\section*{Lời cảm ơn}

Xin cảm ơn thầy Jiang Tao đã đọc bài và đưa ra những ý kiến sửa đổi quý báu.
Xin cảm ơn huấn luyện viên Liu Rujia và bạn Xu Zhilei đã thảo luận sâu với
tác giả về việc viết bài.

\section*{Tài liệu tham khảo}

\begin{enumerate}
  \item Liu Rujia, Huang Liang, \textit{Nghệ thuật thuật toán và thi tin học},
  Nhà xuất bản Đại học Thanh Hoa, 2004.
  \item Xu Zhilei, luận văn đội tuyển tập huấn IOI 2004, \textit{Suffix
  Array}.
  \item CEPC 2003 Problem D, \textit{Dice Contest}.
  \item USACO 2004 February Contest Green Problems.
\end{enumerate}

\appendix
\section{Ghi chú về phụ lục gốc}

Bản gốc kèm phát biểu bài cho ví dụ xúc xắc, RMQ, bài USACO \textit{Distance
Queries}, cùng mã nguồn minh họa bằng C/C++. Theo quy ước của bản dịch này và
theo yêu cầu không dịch code templates, phần mã nguồn không được chuyển ngữ.
Các ý tưởng thuật toán của những chương trình đó đã được trình bày trong thân
bài.

\paragraph{Farewell, My Friend!}
Bài xúc xắc cho sáu trọng số $f(1),\ldots,f(6)$, vị trí đầu
$(x_1,y_1)$ và vị trí cuối $(x_2,y_2)$ trên dải bốn hàng vô hạn theo chiều
ngang. Ban đầu mặt $6$ ở đáy, mặt $1$ ở trên, mặt $2$ hướng ra trước. Cần in
chi phí nhỏ nhất để lăn tới đích.

\paragraph{RMQ tổng quát.}
Cho $N$ số nguyên $A[1..N]$ và $M$ cặp chỉ số $(x,y)$ với $x\le y$. Với mỗi
cặp, in chỉ số của phần tử nhỏ nhất trong $A[x..y]$.

\paragraph{Distance Queries.}
Bài USACO cho một cây đường nông trại với các cạnh ngang dọc có độ dài. Sau
khi đọc mô tả cây, cần trả lời nhiều truy vấn khoảng cách giữa hai nông trại,
tức tổng độ dài các cạnh trên đường đi duy nhất nối chúng.

\end{document}
